\paragraph{Separate resource-oriented gain selection.}
The normal-supply comparison is additionally challenged by a new resource-oriented search, declared after the earlier results but before its own outcomes. Each of scalar hybrid and active SMC receives the same nine $(g_t,g_a)$ candidates, $\{0.7,1.0,1.3\}\times\{0.3,0.5,0.7\}$, and three new development seeds 5600--5602. All 54 development paths use the complete 1800 s repeated-load schedule, fivefold leakage and 1 ms RK4. Holding authority, sensing, observer, 8 s latch and discharge gate are shared. Candidates are ranked lexicographically by number of failed development paths, mean $B$ in Equation~\eqref{eq:boundary-resource}, and candidate ID. The feasibility screen requires confirmed hold, both 10 mm limits and no relief or stroke event. A zero-failure candidate is development-feasible; if none exists the least-failing candidate remains explicitly infeasible. Selection is frozen before eight additional 1800 s validation paths: seeds 5620--5623 under each selected policy. Every candidate and outcome is retained in S13. This is a bounded gain search, not a global resource frontier: SMC's 1200 N boundary-layer amplitude and 0.006 m/s width, hybrid thresholds and other architecture constants remain fixed. Validation concerns normal supply and the nominal plant with the specified leakage; it supplies no new joint-uncertainty or interruption guarantee.
